\setcounterpageref{ex}{0}
\TaoPhieuKiemTra{08}

\setcounter{ex}{0}
\setcounter{bt}{0}

\textbf{\large I. PHẦN TRẮC NGHIỆM (3,0 điểm)}

{\itshape Khoanh tròn vào chữ cái đặt trước câu trả lời đúng. Mỗi câu đúng được 0,5 điểm.}

\begin{ex} % Câu 1
Trên tia số, số đứng ngay liền trước số $40$ là
\choice
{$41$}
{$30$}
{$50$}
{\True $39$}
\loigiai{
Số đứng liền trước số $40$ là $39$.\\
\textbf{Đáp án đúng: D}
}
\end{ex}

\begin{ex} % Câu 2
Biết hai số hạng là $45$ và $4$. Tổng của hai số đó là
\choice
{\True $49$}
{$41$}
{$50$}
{$48$}
\loigiai{
Tổng là: $45 + 4 = 49$.\\
\textbf{Đáp án đúng: A}
}
\end{ex}

\begin{ex} % Câu 3
Túi gạo nặng $10\text{ kg}$, túi ngô nhẹ hơn túi gạo $3\text{ kg}$. Túi ngô cân nặng là
\choice
{$8\text{ kg}$}
{$13\text{ kg}$}
{\True $7\text{ kg}$}
{$6\text{ kg}$}
\loigiai{
Túi ngô nặng: $10 - 3 = 7\text{ kg}$.\\
\textbf{Đáp án đúng: C}
}
\end{ex}

\begin{ex} % Câu 4
Hình ảnh nào dưới đây cho ta hình ảnh của một đường cong?
\choice
{Cạnh của cái thước kẻ}
{\True Cầu vồng trên bầu trời}
{Cạnh của chiếc bàn học}
{Cột cờ thẳng tắp}
\loigiai{
Cầu vồng có dạng uốn lượn cong tự nhiên.\\
\textbf{Đáp án đúng: B}
}
\end{ex}

\begin{ex} % Câu 5
Điền dấu: $84 - 32 \dots 50$.
\choice
{\True $>$}
{$<$}
{$=$}
{Không so sánh được}
\loigiai{
$84 - 32 = 52$. Vì $52 > 50$ nên điền dấu $>$.\\
\textbf{Đáp án đúng: A}
}
\end{ex}

\begin{ex} % Câu 6
Tính: $12 - 6 + 9 = \dots$
\choice
{$14$}
{$16$}
{$13$}
{\True $15$}
\loigiai{
$12 - 6 = 6$; $6 + 9 = 15$.\\
\textbf{Đáp án đúng: D}
}
\end{ex}

\vspace{4mm}
\textbf{\large II. PHẦN TỰ LUẬN (7,0 điểm)}

{\itshape Thí sinh trình bày chi tiết bài làm vào giấy kiểm tra.}

\begin{bt}[1,0 điểm]
Tính nhẩm:
\begin{multicols}{2}
\begin{enumerate}[a)]
  \item $8 + 7$
  \item $7 + 8$
  \item $15 - 5$
  \item $17 - 6$
\end{enumerate}
\end{multicols}
\loigiai{
\begin{itemize}
  \item $8 + 7 = 15$
  \item $7 + 8 = 15$
  \item $15 - 5 = 10$
  \item $17 - 6 = 11$
\end{itemize}
}
\end{bt}

\begin{bt}[1,5 điểm]
Đặt tính rồi tính:
\begin{multicols}{3}
\begin{enumerate}[a)]
  \item $39 + 56$
  \item $91 - 25$
  \item $57 + 40$
\end{enumerate}
\end{multicols}
\loigiai{
\begin{itemize}
  \item $39 + 56 = 95$
  \item $91 - 25 = 66$
  \item $57 + 40 = 97$
\end{itemize}
}
\end{bt}

\begin{bt}[1,0 điểm]
Tính giá trị biểu thức kèm đơn vị đo:
\begin{enumerate}[a)]
  \item $24\text{ kg} + 7\text{ kg} - 11\text{ kg}$
  \item $33\text{ }l - 14\text{ }l + 6\text{ }l$
\end{enumerate}
\loigiai{
\begin{enumerate}[a)]
  \item $24\text{ kg} + 7\text{ kg} - 11\text{ kg} = 31\text{ kg} - 11\text{ kg} = 20\text{ kg}$.
  \item $33\text{ }l - 14\text{ }l + 6\text{ }l = 19\text{ }l + 6\text{ }l = 25\text{ }l$.
\end{enumerate}
}
\end{bt}

\begin{bt}[1,5 điểm]
Bạn Quân có $23$ viên bi. Bạn Sơn có ít hơn bạn Quân $6$ viên bi. Hỏi bạn Sơn có bao nhiêu viên bi?
\loigiai{
Số viên bi bạn Sơn có là:
\[ 23 - 6 = 17 \text{ (viên bi)} \]
\textbf{Đáp số:} $17$ viên bi.
}
\end{bt}

\begin{bt}[1,0 điểm]
Một thùng chứa đầy $43\text{ }l$ dầu ăn. Buổi sáng cửa hàng đã bán ra $12\text{ }l$. Hỏi trong thùng còn lại bao nhiêu lít dầu ăn?
\loigiai{
Trong thùng còn lại số lít dầu ăn là:
\[ 43 - 12 = 31\text{ }l \]
\textbf{Đáp số:} $31\text{ }l$ dầu ăn.
}
\end{bt}

\begin{bt}[1,0 điểm]
Điền số thích hợp vào chỗ chấm: $2, 5, 8, 11, \dots, \dots$
\loigiai{
Quy luật: Dãy số tăng đều $3$ đơn vị.\\
Số tiếp theo: $11 + 3 = 14$.\\
Số tiếp theo nữa: $14 + 3 = 17$.\\
\textbf{Đáp số:} $14$ và $17$.
}
\end{bt}
